% SNAS 2026 double-blind short-paper draft.
% Compile with XeLaTeX so the document uses the installed Times New Roman font.
\documentclass[12pt,letterpaper,twoside]{article}

\usepackage[margin=1in]{geometry}
\usepackage{fontspec}
\setmainfont{Times New Roman}
\usepackage{setspace}
\doublespacing
\usepackage{booktabs}
\usepackage{array}
\usepackage{fancyhdr}
\usepackage{titlesec}
\usepackage{enumitem}
\usepackage{xurl}
\urlstyle{same}
\usepackage[hidelinks]{hyperref}
\hypersetup{
  pdftitle={Beyond Stable Answers: Short-Paper Draft},
  pdfauthor={},
  pdfsubject={Double-blind draft for the 5th SNAS Interdisciplinary Research Conference},
  pdfkeywords={explainable artificial intelligence, vision-language models, construction safety}
}

\setlength{\parindent}{0.5in}
\setlength{\parskip}{0pt}
\setlength{\emergencystretch}{3em}
\raggedbottom
\setlist{nosep,leftmargin=0.75in}

\titleformat{\section}
  {\normalfont\bfseries\centering}
  {\thesection.}{0.5em}{}
\titlespacing*{\section}{0pt}{12pt}{6pt}
\titleformat{\subsection}
  {\normalfont\bfseries}
  {\thesubsection.}{0.5em}{}
\titlespacing*{\subsection}{0pt}{9pt}{3pt}

\pagestyle{fancy}
\fancyhf{}
\fancyhead[LE,LO]{\textbf{DRAFT - NOT CAMERA READY}}
\fancyhead[RE,RO]{\thepage}
\renewcommand{\headrulewidth}{0.4pt}
\setlength{\headheight}{15pt}

\newenvironment{apareferences}
  {\begin{list}{}{%
    \setlength{\leftmargin}{0.5in}%
    \setlength{\itemindent}{-0.5in}%
    \setlength{\itemsep}{0pt}%
    \setlength{\parsep}{0pt}%
    \setlength{\topsep}{0pt}}}
  {\end{list}}

\begin{document}

\begin{center}
\textbf{Beyond Stable Answers: Auditing Explanation Faithfulness in a Construction-Safety Vision-Language Pipeline}
\end{center}

\begin{center}
\textbf{Abstract}
\end{center}

Vision-language models (VLMs) can support construction-safety monitoring by connecting natural-language safety rules to visual evidence. However, answer accuracy alone does not establish whether the evidence underlying a safety judgment is faithful or robust. This feasibility study adapts a six-dimensional explainable artificial intelligence evaluation framework, originally developed for tabular intrusion detection, to a VLM-grounded construction-safety pipeline. We evaluated Florence-2-base-ft on 163 held-out images from ConstructionSite 10k representing compliant scenes, personal protective equipment violations, fall hazards, and struck-by risks. Because Florence-2 does not provide native free-form visual question answering in this implementation, open-vocabulary grounding identified workers and rule-relevant safety objects, while a deterministic person-relative rule produced the compliance judgment. Explanation quality was tested through rule-aware region ranking, safety-object masking, severity-controlled perturbations, targeted and random-location occlusion, size-invariant centroid drift, bounding-box disappearance, worker-loss correction, bootstrap confidence intervals, and paired Wilcoxon tests. Area ranking selected the worker rather than the safety object in 50.3\% of usable cases; rule-aware ranking reduced this to 0.6\%. The corrected rule-aware top-region ablation changed 37.4\% of judgments. Targeted object occlusion changed 39.2\% of judgments and produced greater explanation drift than the strongest random-location control (0.198 versus 0.142; paired \textit{n} = 144, \textit{p} = .0011). Although intersection over union initially classified 25.0\% of stable-answer cases as drifted, only 13.7\% moved more than 0.2 of the image diagonal, revealing substantial small-box bias; 5.1\% of targeted explanations disappeared entirely. These findings show that trustworthy construction-safety AI requires explanation-level evaluation, semantic region selection, size-aware metrics, and explicit failure accounting in addition to answer-level performance. The pipeline remains a research testbed requiring human oversight, not an autonomous safety system.

\noindent\textit{Keywords:} explainable artificial intelligence; vision-language models; construction safety; explanation faithfulness; robustness; trustworthy AI

\section{Introduction}

Construction remains a high-risk industry. In 2024, construction and extraction occupations experienced 1,032 fatal work injuries in the United States, or 12.6 deaths per 100,000 full-time-equivalent workers (U.S. Bureau of Labor Statistics [BLS], 2026). Manual inspection remains essential, but visual monitoring systems may help professionals review site imagery and prioritize follow-up investigation.

VLMs combine image inputs with natural-language instructions and have been studied for construction hazard classification, visual grounding, and caption generation (Chen \& Zou, 2025; Kim et al., 2025; Li et al., 2026). Yet task accuracy does not establish that a model used relevant evidence. A safety answer may remain stable while its grounding moves to an unrelated object, and a plausible rationale may not reflect the process that generated an output (Jain \& Wallace, 2019; Li et al., 2025). These problems are consequential in construction because an apparently stable warning is unsuitable for professional review if its evidence relocates or disappears.

This study adapts the descriptive accuracy, sparsity, stability, efficiency, robustness, and completeness dimensions of an XAI evaluation framework originally developed for network intrusion detection (Arreche et al., 2024; Arreche \& Abdallah, 2025). The dimensions transfer conceptually, but image regions introduce ranking choices, box-size effects, masking artifacts, and detector fallback paths. The study asks: (1) how often do safety answers remain unchanged while visual evidence relocates or disappears; (2) how much apparent explanation drift is caused by size-sensitive intersection over union (IoU) or inappropriate region ranking; and (3) whether occluding a rule-relevant object produces greater change than random-location occlusion. The contribution is a construct-validity audit that separates output change, evidence movement, and evidence availability.

\section{Related Work}

Recent construction research demonstrates the promise of contextual visual models. ConstructionSite 10k contains 10,013 images with captions, safety-rule questions, rationales, and grounding annotations (Chen \& Zou, 2025). Other work has evaluated general-purpose VLMs, adapted VLMs for fall hazards, and combined object detectors with smaller VLMs to improve hazard recognition and explanation quality (Adil et al., 2025, 2026; Kim et al., 2025; Li et al., 2026). These studies primarily establish task capability; they do not fully test whether the visual evidence underlying a judgment is stable and causally relevant.

Faithfulness concerns whether an explanation reflects the reason for an output or the process that produced it (Phillips et al., 2021). Attention should not automatically be interpreted as explanation, and perturbation metrics can reward weak attributions when their assumptions are not validated (Jain \& Wallace, 2019; Wu et al., 2024). In workplace settings, the NIST AI Risk Management Framework and NIOSH guidance emphasize context, documentation, human oversight, and risk controls (Howard \& Schulte, 2024; National Institute of Standards and Technology [NIST], 2023). Accordingly, this paper evaluates technical faithfulness rather than claiming that explanations improve human trust or injury outcomes.

\section{Method}

\subsection{Data and Model}

The sample contained 163 images drawn exclusively from the held-out test split of ConstructionSite 10k: 50 compliant scenes, 50 personal protective equipment violations, 50 fall hazards, and 13 struck-by risks. The struck-by subgroup was retained but treated as exploratory because it contained fewer than 30 images. Images with broken paths, unreadable files, or invalid dimensions were excluded before inference.

The model was \textit{microsoft/Florence-2-base-ft}, an approximately 231.6-million-parameter encoder-decoder model supporting prompted detection and grounding (Xiao et al., 2024). Experiments ran on one NVIDIA RTX 3070. This Florence-2 implementation did not provide a native free-form visual-question-answering task. It instead grounded a worker and a rule-relevant object---hard hat, harness, guardrail, or excavator---and deterministic person-relative geometry produced the compliant/violation judgment. The evaluated system is therefore a VLM-grounded safety decision pipeline, not native end-to-end VQA.

The first scene-level object-existence proxy detected 3.0\% of 100 PPE and fall violations with 97.4\% specificity on corresponding compliant cases. Person-relative logic improved sensitivity to 27.0\% and produced 73.7\% specificity. This redesign avoided an almost constant compliant classifier but remained inadequate for deployment.

\subsection{Explanation Interventions}

The primary explanation was the model-provided box for the queried safety object. Four images without a native box used a labeled grid fallback. The frozen pilot ranked boxes by pixel area, which often promoted a worker body over small PPE. A corrected rule-aware policy ranked the queried object first and used area only within semantic tiers.

Descriptive accuracy was tested by masking the top-ranked region and rerunning the pipeline. Black, blur, and inpaint operators assessed mask sensitivity. Targeted occlusion masked the baseline safety-object box, while centered and seeded random-location occlusions served as controls across three severities. The robustness sweep contained 2,603 perturbation rows.

Answer change captured output sensitivity. Explanation change used IoU, centroid displacement normalized by image diagonal, and object disappearance. Disappearance was recorded explicitly rather than omitted as missing IoU. Because masking could erase the only detected worker and activate a scene-level fallback, flips co-occurring with worker disappearance were labeled as worker-loss failures and removed from corrected flip rates.

\subsection{Statistical Analysis}

Headline rates used fixed-seed percentile-bootstrap 95\% confidence intervals. Paired Wilcoxon signed-rank tests compared per-image explanation drift. Targeted occlusion was paired with the strongest random-location condition, severity 0.35, on images containing both measurements, yielding \textit{n} = 144. The frozen outputs were preserved while corrected modes wrote separate artifacts. All eleven pipeline stages were also rerun from a clean environment on a fresh 20-image sample to check reproducibility.

\section{Results}

\subsection{Ranking and Intervention Validity}

All six XAI dimensions ran end to end. Frozen top-region descriptive accuracy was 36.2\% (bootstrap 95\% CI [28.8\%, 43.6\%]); stability reruns produced 77.5\% answer agreement, and four low-severity perturbations produced 80.7\% answer survival. These values establish computational feasibility, not operational safety.

Among 159 images with model-provided boxes, area ranking selected the worker as the top region in 50.3\% of cases. Rule-aware ranking reduced this to 0.6\%. As Table~\ref{tab:ranking} shows, the corrected policy modestly increased the flip rate but, more importantly, changed the construct being tested from large-object dependence to rule-relevant evidence dependence.

\begin{table}[htbp]
\centering
\caption{Top-Region Ranking and Answer-Change Results}
\label{tab:ranking}
\begin{tabular}{@{}lccc@{}}
\toprule
Policy & Worker top-1 & Raw flip & Corrected flip \\
\midrule
Area baseline & 50.3\% & 36.2\% & 34.4\% \\
Rule-aware & 0.6\% & 38.7\% & 37.4\% \\
\bottomrule
\end{tabular}
\end{table}

\subsection{Targeted Versus Control Occlusion}

Targeted occlusion was available for 158 images and changed 39.2\% of judgments. Random-location occlusion changed 17.2\% when pooled across severities and 24.5\% at severity 0.35. In the paired analysis, targeted occlusion produced mean normalized centroid drift of 0.198 compared with 0.142 for the strongest random control (\textit{n} = 144, Wilcoxon \textit{p} = .0011). Table~\ref{tab:occlusion} also shows that explanation disappearance occurred under both targeted and control interventions, supporting separate reporting of evidence availability.

\begin{table}[htbp]
\centering
\caption{Occlusion Outcomes}
\label{tab:occlusion}
\begin{tabular}{@{}lccc@{}}
\toprule
Condition & Answer change & Centroid drift & Disappearance \\
\midrule
Targeted object & 39.2\% & 0.196 overall & 5.1\% \\
Random, severity 0.35 & 24.5\% & 0.143 overall & 6.7\% \\
Center, severity 0.20 & 28.2\% & 0.105 overall & 4.9\% \\
\bottomrule
\end{tabular}
\end{table}

\subsection{Stable Answers and Failure Accounting}

The fixed pilot initially found that 28.3\% of stable-answer rows with comparable boxes had IoU below 0.3. Validity analysis reproduced an IoU-based rate near 23--25\% but showed strong size bias: 31.0\% of small boxes crossed the threshold, compared with 24.4\% of medium and 13.5\% of large boxes. Mean normalized centroid drift was only 0.062, and 13.7\% moved more than 0.2 of the image diagonal. Genuine hidden relocation therefore persisted, but IoU alone overstated its frequency. Targeted explanations disappeared in another 5.1\% of reruns and would be missed by an IoU-only average.

Pipeline tracing also changed the interpretation of answer flips. Under centered occlusion at severity 0.20, the raw answer-change rate was 28.2\%, whereas the worker-loss-corrected rate was 22.7\%; nine of 163 flips coincided with worker loss. For rule-aware top-region masking, correction reduced the rate from 38.7\% to 37.4\%. A perturbation audit must therefore record intermediate detector state so pipeline failure is not misclassified as explanation faithfulness.

\section{Discussion and Implications}

The results demonstrate that answer robustness and explanation robustness are different properties. A stable answer can coexist with relocated or unavailable evidence, while an answer flip can arise from detector failure rather than removal of the intended safety cue. Trustworthy evaluation should report output change, evidence movement, and evidence availability separately.

Three validity controls were consequential. Rule-aware ranking aligned the intervention with the safety concept rather than pixel area. Targeted and random-location occlusion separated object removal from generic disruption. Centroid drift and disappearance prevented IoU from serving as the sole stability measure. These controls show why tabular XAI metrics cannot be transferred to visual evidence without testing the assumptions embedded in region selection and perturbation.

The practical implication is procedural, not a claim of autonomous inspection. A future decision-support tool should expose the visual evidence associated with a warning, disclose when evidence is absent or unstable, and defer to trained safety personnel. Explanations should help professionals interrogate system limits rather than encourage unconditional acceptance. Human-centered studies would still be required to test reliance, workload, and decision quality (Naiseh et al., 2023).

\section{Limitations and Ethical Considerations}

This feasibility study evaluates one small VLM, one dataset, and a modest sample. The struck-by subgroup contains only 13 images, and results may not generalize to native-VQA models, video, or real sites. The final judgment comes from geometry over VLM grounding rather than end-to-end generative reasoning. Open-vocabulary grounding often returned one worker in multi-worker scenes, and the resulting 27.0\% violation sensitivity is unacceptable for deployment. Token-probability confidence was not calibrated and was not interpreted as a correctness probability. Perturbations can create unnatural images, although multiple mask operators, severity sweeps, random controls, and worker-loss accounting reduce that concern.

No human-subject study was conducted. The study does not establish meaningfulness, calibrated trust, workload improvement, or injury reduction. Workplace imagery also creates privacy, surveillance, and fairness risks. The work does not infer worker identity, intent, competence, or blame, and the dataset cannot support demographic fairness claims. Any future use should include data minimization, notice, access controls, human review, appeal procedures, and a prohibition on autonomous disciplinary action.

\section{Conclusion}

A stable construction-safety answer does not guarantee faithful visual evidence. In this 163-image study, semantic ranking, IoU size bias, and worker-detector fallback each changed the interpretation of XAI metrics. Rule-aware ranking aligned interventions with safety objects, targeted occlusion caused greater drift than random controls, and size-invariant analysis showed that hidden evidence relocation persisted at a lower rate than IoU suggested. Trustworthy construction-safety AI requires explanation-level interventions, size-aware metrics, explicit disappearance and fallback accounting, and human oversight. The current pipeline is a reproducible research testbed, not an autonomous safety system.

\label{bodyend}
\clearpage
\section*{References}
\addcontentsline{toc}{section}{References}

\begin{apareferences}
\item Adil, M., Ahmed, M., Aqib, M., Gonzalez, V. A., Lee, G., \& Mei, Q. (2026). Integration of object detection and small VLMs for construction safety hazard identification. \textit{arXiv}. \url{https://arxiv.org/abs/2604.05210}

\item Adil, M., Lee, G., Gonzalez, V. A., \& Mei, Q. (2025). Using vision language models for safety hazard identification in construction. \textit{arXiv}. \url{https://arxiv.org/abs/2504.09083}

\item Arreche, O., \& Abdallah, M. (2025). A comparative analysis of DNN-based white-box explainable AI methods in network security. \textit{EURASIP Journal on Information Security, 2025}, Article 16. \url{https://doi.org/10.1186/s13635-025-00201-x}

\item Arreche, O., Guntur, T. R., Roberts, J. W., \& Abdallah, M. (2024). E-XAI: Evaluating black-box explainable AI frameworks for network intrusion detection. \textit{IEEE Access, 12}, 23954--23988. \url{https://doi.org/10.1109/ACCESS.2024.3365140}

\item Chen, X., \& Zou, Z. (2025). Are large pre-trained vision language models effective construction safety inspectors? \textit{arXiv}. \url{https://arxiv.org/abs/2508.11011}

\item Howard, J., \& Schulte, P. A. (2024, September 9). Exploring approaches to keep an AI-enabled workplace safe for workers. \textit{National Institute for Occupational Safety and Health}. \url{https://www.cdc.gov/niosh/bulletin/2024/ai-risk-management.html}

\item Jain, S., \& Wallace, B. C. (2019). Attention is not explanation. In \textit{Proceedings of NAACL-HLT 2019} (pp. 3543--3556). \url{https://doi.org/10.18653/v1/N19-1357}

\item Kim, T., Kim, S., Chern, W.-C., Park, S., Kim, D., \& Kim, H. (2025). Optimizing large vision-language models for context-aware construction safety assessment. \textit{Automation in Construction, 180}, 106510. \url{https://doi.org/10.1016/j.autcon.2025.106510}

\item Li, K., Vosselman, G., \& Yang, M. Y. (2025). Multimodal rationales for explainable visual question answering. In \textit{Proceedings of the CVPR Workshops} (pp. 191--201). \url{https://openaccess.thecvf.com/content/CVPR2025W/MULA2025/html/Li_Multimodal_Rationales_for_Explainable_Visual_Question_Answering_CVPRW_2025_paper.html}

\item Li, Y., Xu, F., Zhang, Z., Mei, X., \& Huang, H. (2026). Construction site fall hazard identification and automated captioning using adapted vision-language models. \textit{Automation in Construction, 183}, 106790. \url{https://doi.org/10.1016/j.autcon.2026.106790}

\item Naiseh, M., Al-Thani, D., Jiang, N., \& Ali, R. (2023). How the different explanation classes impact trust calibration: The case of clinical decision support systems. \textit{International Journal of Human-Computer Studies, 169}, 102941. \url{https://doi.org/10.1016/j.ijhcs.2022.102941}

\item National Institute of Standards and Technology. (2023). \textit{Artificial Intelligence Risk Management Framework (AI RMF 1.0)}. \url{https://doi.org/10.6028/NIST.AI.100-1}

\item Phillips, P. J., Hahn, C. A., Fontana, P. C., Yates, A. N., Greene, K., Broniatowski, D. A., \& Przybocki, M. A. (2021). \textit{Four principles of explainable artificial intelligence} (NISTIR 8312). National Institute of Standards and Technology. \url{https://doi.org/10.6028/NIST.IR.8312}

\item U.S. Bureau of Labor Statistics. (2026). \textit{Number and rate of fatal work injuries, civilian workers, by major occupational group, 2024}. \url{https://www.bls.gov/charts/census-of-fatal-occupational-injuries/number-and-rate-of-fatal-work-injuries-by-occupation.htm}

\item Wu, J., Kang, W., Tang, H., Hong, Y., \& Yan, Y. (2024). On the faithfulness of Vision Transformer explanations. In \textit{Proceedings of the IEEE/CVF Conference on Computer Vision and Pattern Recognition} (pp. 10936--10945). \url{https://openaccess.thecvf.com/content/CVPR2024/html/Wu_On_the_Faithfulness_of_Vision_Transformer_Explanations_CVPR_2024_paper.html}

\item Xiao, B., Wu, H., Xu, W., Dai, X., Hu, H., Lu, Y., Zeng, M., Liu, C., \& Yuan, L. (2024). Florence-2: Advancing a unified representation for a variety of vision tasks. In \textit{Proceedings of the IEEE/CVF Conference on Computer Vision and Pattern Recognition}. \url{https://arxiv.org/abs/2311.06242}
\end{apareferences}

\end{document}
